% Tapis de Quillen, with Quillen's offprint and preprint (cote 162-5, typed part): archive cote 162-5, site page 34 (PDF page 35).
% typescript. Draft compared: none.
% AI-transcribed by Claude Opus 5.5 (claude-opus-5-5, adaptive thinking, effort high), Message Batches API, batch t162_full (msgbatch_01UAh3zyWPd454AeDkhnmQ33), 2026-11-13; not yet checked by eye.
\scanpage{162-5/34}
premier.  Comme applic\xout{t}ation, Quillen obtient une interprétation faisceauti-%
que de la cohomologie d'un ensemble semi-simplicial à coéfficients dans un
système \inserted{local} con\xout{tra}variant \inserted{F} \xout{\unsure{dx}}(défini classiquement par le complexe cosimplicial
\margin{$\prod$ ?} des  $C^n(F) = \coprod_{x\in X_n} F(x)$ ):  on considère le système local contravariant
défini par F, on l'interprète comme un faisceau sur \inserted[typed]{T(X) i.e. objet de} $(\mathrm{Ssimpl})_{/X}$ , et on
prend sa cohomologie. - Cependant, quand F est un système de coéfficien\unsure{ts|ss}
co\unsure{va}riant pas nécessairement local, on n'a tjrs pas d'interp\xout{\unsure{r}}étation de se\unsure{s|n}
groupes de cohomologie  classiques en t ermes faisc\xout{\unsure{e}}autiques; ni, lorsque F
est \unsure{contravariant}, de son homologie, ou inversement de sa cohomologie
faisceautique en termes classiques.

A propos de la notion de foncteur qui est un homotopisme. Quillen
montre qu\xout{e}'un tel foncteur $F: C \hand{\longrightarrow} C'$ induit une équiva\unsure{lence} entre la sous-%
catégorie triangulée \inserted{$D^b_{lc}(C')$} de la catégorie \inserted{dérivée \inserted{bornée} de celle \deleted{des}} des faisceaux abéliens  sur C', dont
les faisceaux de cohomologie sont des systêmes locaux, et la catégorie ana-%
logue pour C; \xout{\unsure{Ilxsxxx}}  et réciproquement. \xout{Pourxlaxpartiexdirecte} On peut
dans cet énoncé introduire aussi n'importe quel anneau de base (à condition
de le supposer $\neq$ 0 dans le cas de la réciproque); \xout{\unsure{constant}} la partie directe
vaut aussi avec un anneau de boefficients pas néc. constant, mais constant
tordu.  Je pense que ce résultat (facile) doit pouvoir se généraliser ainsi:
Soit  f: X \quad X' un morphisme de topos qui  soit tel que pour tout \xout{système}
faisceau loc. const. sur X', f induise un isomorphisme sur les cohomologies
(avec cas non commutatif inclus).  Supposons que \deleted{\xout{\unsure{le topos}} X \unsure{soit}} \inserted{X et $X'$} \uline{localement}
\margin{\unsure{$\circ$}} \margin{$\|$} \uline{homotopiquement trivial}, i.e. que \xout{tout} pour tout entier n\inserted{$\geq 1$}, tout objet U \deleted{de}
\deleted{$X'$} ait un \inserted[typed]{recouvrement} \xout{recouvrement} par des $U_i \hand{\longrightarrow} U$, tels que  a) tout système local sur U
devient constant sur $U_i$ , et toute section sur U devient constante sur $U_i$ ,
et  b)  pour tout \xout{système}  groupe abélien G, les $H^{\hand{j}}(U,G) \longrightarrow H^j(U_i,G)$ sont
nuls pour $1 \hand{\leq} j \hand{\leq} n$.  Alors le foncteur $D^b_{lc}(X') \longrightarrow D^b_{lc}(X)$  induit par f est

\margin{\unsure{$\circ$} Attention, cette condition n'est \unsure{pratiquement} \uline{pas} satisfaite par les schémas avec leur topologie étale (mais bien par \unsure{schémas loc. noeth.} avec top. Zariski).}

\mspage{34}
